% GENERATED FILE. Edit canonical lessons, then run npm run generate:books.
% canonical-source-hash: fnv1a64:66aebb9e1ce38a33
% canonical-lessons: AR-C425-mudif, AR-C425-rakib, AR-C425-wijha, AR-C425-shadda, AR-C425-wusul

\chapter{On board}
\label{ch:ar-on-board}
\hlchaptermodality{425}

\begin{chapteropening}
\textbf{By the end of this chapter:} \emph{I can say a flight attendant, a passenger, a destination, to tighten and arrival.}

Complete the last lesson of chapter 425: I can say a flight attendant, a passenger, a destination, to tighten and arrival.
\end{chapteropening}

\section[\texorpdfstring{\ar{مضيف} (muḍīf)}{مضيف (muḍīf)}]{\ar{مضيف} (muḍīf) — a flight attendant; a host}

\label{lesson:AR-C425-mudif}



\begin{quote}
Before the new one: say the Arabic for a seat belt, then the Arabic for a crew.
\end{quote}

\subsection*{You'll want to know: \ar{مضيف}}
\textbf{\ar{مضيف}} — \emph{muḍīf} — \textquotedblleft{}a flight attendant; a host\textquotedblright{}.

\begin{quote}
\ar{أحضر} \ar{المضيف} \ar{كوب} \ar{ماء} \ar{لأمي}.
\end{quote}

\begin{quote}
\emph{aḥḍara al-muḍīf kūb māʾ li-ummī.}
\end{quote}

\begin{quote}
The flight attendant brought my mother a cup of water.
\end{quote}

A woman is \emph{muḍīfa}.

\begin{cousinweb}
\textbf{\ar{مضيف}} (\emph{muḍīf}) is built on the native Arabic root \emph{ḍ-y-f}, \textquotedblleft{}to be a guest\textquotedblright{} (Lane). It is the Form IV participle, \textquotedblleft{}one who takes in guests\textquotedblright{}: a host, and on a plane a flight attendant.
\end{cousinweb}

\begin{grammarlens}[title={What you've built}]
One more word for getting around.
\end{grammarlens}

\subsection*{Your turn}
\begin{itemize}
\raggedright
  \item \emph{Say it:} \emph{muḍīf}
  \item \emph{Say it:} \emph{muḍīf}, once more
  \item \emph{Say it:} \emph{aḥḍara al-muḍīf kūb māʾ li-ummī.} — The flight attendant brought my mother a cup of water.
  \item \emph{Recall:} say the Arabic for a seat belt, then the Arabic for a crew, then say \emph{muḍīf} again
\end{itemize}

\begin{tcolorbox}[breakable,colback=teal!4,colframe=teal!35!black,title={Before you move on}]
What does \ar{مضيف} mean? (\textquotedblleft{}A flight attendant; a host\textquotedblright{}.) Say it once more.
\end{tcolorbox}
\section[\texorpdfstring{\ar{راكب} (rākib)}{راكب (rākib)}]{\ar{راكب} (rākib) — a passenger; a rider}

\label{lesson:AR-C425-rakib}



\begin{quote}
Before the new one: say the Arabic for a crew, then the Arabic for a flight attendant.
\end{quote}

\subsection*{You'll want to know: \ar{راكب}}
\textbf{\ar{راكب}} — \emph{rākib} — \textquotedblleft{}a passenger; a rider\textquotedblright{}.

\begin{quote}
\ar{نام} \ar{راكب} \ar{بجانبي} \ar{أثناء} \ar{الرحلة}.
\end{quote}

\begin{quote}
\emph{nāma rākib bi-jānibī athnāʾa ar-riḥla.}
\end{quote}

\begin{quote}
A passenger beside me slept during the flight.
\end{quote}

Plural: \ar{ركاب}, \emph{rukkāb}.

\begin{cousinweb}
\textbf{\ar{راكب}} (\emph{rākib}) is built on the native Arabic root \emph{r-k-b}, \textquotedblleft{}to ride\textquotedblright{} (Lane). It is the active participle of \emph{rakiba}, \textquotedblleft{}to ride\textquotedblright{}: one who rides.
\end{cousinweb}

\begin{grammarlens}[title={What you've built}]
One more word for getting around.
\end{grammarlens}

\subsection*{Your turn}
\begin{itemize}
\raggedright
  \item \emph{Say it:} \emph{rākib}
  \item \emph{Say it:} \emph{rākib}, once more
  \item \emph{Say it:} \emph{nāma rākib bi-jānibī athnāʾa ar-riḥla.} — A passenger beside me slept during the flight.
  \item \emph{Recall:} say the Arabic for a crew, then the Arabic for a flight attendant, then say \emph{rākib} again
\end{itemize}

\begin{tcolorbox}[breakable,colback=teal!4,colframe=teal!35!black,title={Before you move on}]
What does \ar{راكب} mean? (\textquotedblleft{}A passenger; a rider\textquotedblright{}.) Say it once more.
\end{tcolorbox}
\section[\texorpdfstring{\ar{وجهة} (wijha)}{وجهة (wijha)}]{\ar{وجهة} (wijha) — a destination; a direction}

\label{lesson:AR-C425-wijha}



\begin{quote}
Before the new one: say the Arabic for a flight attendant, then the Arabic for a passenger.
\end{quote}

\subsection*{You'll want to know: \ar{وجهة}}
\textbf{\ar{وجهة}} — \emph{wijha} — \textquotedblleft{}a destination; a direction\textquotedblright{}.

\begin{quote}
\ar{وجهتنا} \ar{مدينة} \ar{على} \ar{البحر}.
\end{quote}

\begin{quote}
\emph{wijhatunā madīna ʿalā al-baḥr.}
\end{quote}

\begin{quote}
Our destination is a city by the sea.
\end{quote}

\begin{cousinweb}
\textbf{\ar{وجهة}} (\emph{wijha}) is built on the native Arabic root \emph{w-j-h}, \textquotedblleft{}the face; to face\textquotedblright{} (Lane). A \emph{wijha} is the place one turns one's face towards.
\end{cousinweb}

\begin{grammarlens}[title={What you've built}]
One more word for getting around.
\end{grammarlens}

\subsection*{Your turn}
\begin{itemize}
\raggedright
  \item \emph{Say it:} \emph{wijha}
  \item \emph{Say it:} \emph{wijha}, once more
  \item \emph{Say it:} \emph{wijhatunā madīna ʿalā al-baḥr.} — Our destination is a city by the sea.
  \item \emph{Recall:} say the Arabic for a flight attendant, then the Arabic for a passenger, then say \emph{wijha} again
\end{itemize}

\begin{tcolorbox}[breakable,colback=teal!4,colframe=teal!35!black,title={Before you move on}]
What does \ar{وجهة} mean? (\textquotedblleft{}A destination; a direction\textquotedblright{}.) Say it once more.
\end{tcolorbox}
\section[\texorpdfstring{\ar{شد} (shadda)}{شد (shadda)}]{\ar{شد} (shadda) — to tighten; to fasten}

\label{lesson:AR-C425-shadda}



\begin{quote}
Before the new one: say the Arabic for a passenger, then the Arabic for a destination.
\end{quote}

\subsection*{You'll want to know: \ar{شد}}
\textbf{\ar{شد}} — \emph{shadda} — \textquotedblleft{}to tighten; to fasten\textquotedblright{}.

\begin{quote}
\ar{شد} \ar{الراكب} \ar{حزامه} \ar{قبل} \ar{الهبوط}.
\end{quote}

\begin{quote}
\emph{shadda ar-rākib ḥizāmahu qabla al-hubūṭ.}
\end{quote}

\begin{quote}
The passenger fastened his belt before landing.
\end{quote}

Present: \ar{يشد}, \emph{yashudd}.

A doubled verb: \emph{shadadtu}, \textquotedblleft{}I tightened\textquotedblright{}, unfolds the two d's.

\begin{cousinweb}
\textbf{\ar{شد}} (\emph{shadda}) is built on the native Arabic root \emph{sh-d-d}, \textquotedblleft{}to be firm; to bind tightly\textquotedblright{} (Lane). It is the plain verb of a doubled root.
\end{cousinweb}

\begin{grammarlens}[title={What you've built}]
One more word for getting around.
\end{grammarlens}

\subsection*{Your turn}
\begin{itemize}
\raggedright
  \item \emph{Say it:} \emph{shadda}
  \item \emph{Say it:} \emph{shadda}, once more
  \item \emph{Say it:} \emph{shadda ar-rākib ḥizāmahu qabla al-hubūṭ.} — The passenger fastened his belt before landing.
  \item \emph{Recall:} say the Arabic for a passenger, then the Arabic for a destination, then say \emph{shadda} again
\end{itemize}

\begin{tcolorbox}[breakable,colback=teal!4,colframe=teal!35!black,title={Before you move on}]
What does \ar{شد} mean? (\textquotedblleft{}To tighten; to fasten\textquotedblright{}.) Say it once more.
\end{tcolorbox}
\section[\texorpdfstring{\ar{وصول} (wuṣūl)}{وصول (wuṣūl)}]{\ar{وصول} (wuṣūl) — arrival}

\label{lesson:AR-C425-wusul}



\begin{quote}
Before the new one: say the Arabic for a destination, then the Arabic for to tighten.
\end{quote}

\subsection*{You'll want to know: \ar{وصول}}
\textbf{\ar{وصول}} — \emph{wuṣūl} — \textquotedblleft{}arrival\textquotedblright{}.

\begin{quote}
\ar{اتصلت} \ar{بأمي} \ar{بعد} \ar{الوصول}.
\end{quote}

\begin{quote}
\emph{ittaṣaltu bi-ummī baʿda al-wuṣūl.}
\end{quote}

\begin{quote}
I phoned my mother after arrival.
\end{quote}

\begin{cousinweb}
\textbf{\ar{وصول}} (\emph{wuṣūl}) is built on the native Arabic root \emph{w-ṣ-l}, \textquotedblleft{}to arrive, to reach\textquotedblright{} (Lane). It is the verbal noun of \emph{waṣala}, \textquotedblleft{}to arrive\textquotedblright{}.
\end{cousinweb}

\begin{grammarlens}[title={What you've built}]
One more word for getting around.
\end{grammarlens}

\subsection*{Your turn}
\begin{itemize}
\raggedright
  \item \emph{Say it:} \emph{wuṣūl}
  \item \emph{Say it:} \emph{wuṣūl}, once more
  \item \emph{Say it:} \emph{ittaṣaltu bi-ummī baʿda al-wuṣūl.} — I phoned my mother after arrival.
  \item \emph{Recall:} say the Arabic for a destination, then the Arabic for to tighten, then say \emph{wuṣūl} again
\end{itemize}

\begin{tcolorbox}[breakable,colback=teal!4,colframe=teal!35!black,title={Before you move on}]
What does \ar{وصول} mean? (\textquotedblleft{}Arrival\textquotedblright{}.) Say it once more.
\end{tcolorbox}
